% AI-written, AI cross-reviewed; no human sign-off; not a claim that RH is proved
% Atlas unit P029. Local manuscript v3, item 7.6, heading on printed PDF p. 75.
% UpperTheory.GalN.V.Win.N.countK_alt
% paper2/inputs_one.tex:242–269. Verbatim excerpt; SHA-256 85ffc664734e43e3a1243ca745c6409fe505eb2b0c8cc5b9aaf0d7ea5a6c9562.
% Original TeX uses the macros and cross-reference keys of the source manuscript.


% Start of atlas edition annotation. The following block is outside the original excerpt.
\paragraph{Status of named inputs in this reading.}
\noindent\nolinkurl{UpperTheory.GalN.V.Win.N.countK_alt} --- in-repo constant (\path{Development/AutomorphicArithmetic/AutomorphicVarieties/ModularityArgument/ThresholdArguments/DenseThresholdCases/FrobeniusTraceCountBoundOfCmEigensystemCount.lean:63--125}); not compared.\par
% End of atlas edition annotation.

\begin{formalstatement}[A CM classifier of fixed purity]\label{pn:P029}
Assume the automorphic hypotheses \(H'\) and the CM level and
Satake--Frobenius correspondence \(\mathrm{LevelKNAlt}\).
For each CM field \(K\) and \(n\ge2\) there exist \(c,A\in\mathbb N\).
Let \(m_1>\cdots>m_n\) satisfy \(m_i+m_{n+1-i}=1-n\)
and \(|m_i|\le H\). Given \(M\ge1\) and \(w\mid p\) with \(p\) prime,
\(p\nmid M\operatorname{disc}(K)\), there are eigenvalue data \(D\)
at \(Nw\) of weight \(m\) and purity \(1-n\), with
\[
 |D.S|\le c[M(H+1)]^A .
\]
For every odd prime \(\ell\ne p\), closure and coefficient comparison
\(\iota\), these same data classify the admitted
\(\mathrm{AdmKNAlt}(\iota,M,m)\) representations of \(G_K\) at \(w\).
\end{formalstatement}\nodestatus{P029}

The written argument defines the tuple \(v\mapsto (Nv)^{(n-1)/2}B_v\)
of normalized Satake multisets at the places \(v\nmid M\) on isomorphism
classes. Barnet-Lamb, Gee, Geraghty and Taylor, \emph{Potential automorphy
and change of weight}, Theorem~2.1.1 and the end of \S5.1, give
compatible Frobenius data for each regular algebraic polarized cuspidal
representation. To obtain one finite classifier the source uses the
cohomological bound on the number of cusp representations, conjugation
of their full Satake tuples, integrality in the prescribed normalization,
and separation of semisimple representations by almost-everywhere
Frobenius data. The simultaneous integral comparison with the polynomial
level and weight bound is an adaptation; its full proof is in Lean only
(Blueprint Z1.5.03~\bploc{r8:small:Z1.5.03}). \cite{r21-blggt}
